\scanpage{14}
Une normalisation \deleted{\unsure{est}} un peu plus poussée est obtenue ainsi : On se donne une trivialisation de $\mathcal L$ par une section partout non nulle $D$ de $\mathcal L$, [dérivation invariante], on choisit alors la base $1,x,y$, de $\mathcal E_3$ de façon que l'image de $y$ dans $\mathcal L_2$ soit $i_2(D)$, celle de $x$ dans $\mathcal L_3$ soit $i_3(D)$. \deleted{\unsure{Celà|Cet}}

\deleted{$[$} i.e. on prend comme structure auxiliaire la donnée d'un $\mathcal L$ \inserted{inv.}, d'un $\mathcal E$ loc. libre de rang 3, muni d'un \unsure{splitting|splittage} où les facteurs successifs sont munis un isom avec $\mathcal L^0,\mathcal L^2,\mathcal L^3$ respectivem. Le groupe d'automorphismes de cette structure est \deleted{formé} \unsure{le|donc le} sous-groupe de $G_m\times\mathrm{Gl}(3)$, formé des \deleted{\unsure{triplets}} \inserted{couples} de la forme
\[
\begin{aligned}
\left(\lambda,\begin{pmatrix}1&\alpha&\beta\\0&\lambda^2&\gamma\\0&0&\lambda^3\end{pmatrix}\right)
&\qquad \lambda\in\Gamma(G_{m_S}/S)=\Gamma(S,\underline O_S^{\unsure{*|\times}})\\
&\qquad \alpha,\beta,\gamma\in\Gamma(S,\underline O_S)
\end{aligned}
\]

\deleted{La fixation de la structure ainsi conduit à une équation} Isomorphe \deleted{\unsure{i.e.}} au sous-groupe du \deleted{\unsure{précédent}} \inserted{groupe plus haut}
\[
\begin{pmatrix}1&\alpha&\beta\\0&\lambda&\gamma\\0&0&\mu\end{pmatrix}
\quad\text{pour}\quad \lambda=\left(\frac{\mu}{\lambda}\right)^2,\quad \mu=\left(\frac{\mu}{\lambda}\right)^3
\quad\text{i.e.}\quad \boxed{\mu^2-\lambda^3=0}.
\]

L'équation de $C$ pour une telle base \deleted{prend alors la forme} [caractérisée par la condition $\underline{x^2-y^3\in\mathcal E_5}$] prend la forme
\[
x^2-y^3+bxy+cy^2+dx+ey+f=0
\]
[i.e. on prend $a=-1$ en l'équation plus haut]
\[
\boxed{\delta(b,c,d,e,f)\text{ une unité}}
\]
% archive sheet number: 13
